\documentclass[../main]{subfiles}
\begin{document}
\chapter{Teorema de Bernoulli}
\img{images/03/bernoulli.jpg}{Análisis aerodinámico sobre Cybertruck Tesla Motors inc.}

\info{
Contenido}{
Teorema de Bernoulli, teorema de Torricelli
}

\section{Teorema de Bernoulli}

\begin{center}
\includegraphics[scale=0.5]{images/03/bernoulli.png}
\end{center}

\begin{equation}
     \dfrac{p_1}{\gamma}+\dfrac{v_1^2}{2g}+h_1=\dfrac{p_2}{\gamma}+\dfrac{v_2^2}{2g}+h_2
\end{equation}

\info{Teorema de Bernoulli}{

El teorema de Bernoulli establece que en un fluido ideal en movimiento, la suma de la presión estática, la presión dinámica y la energía potencial gravitatoria por unidad de masa es constante a lo largo de una corriente de fluido.

En su forma matemática, el teorema de Bernoulli se expresa de la siguiente manera:



\[ P + \frac{1}{2} \rho v^2 + \rho g h = \text{constante} \]

Donde:
- \( P \) es la presión estática del fluido.

- \( \rho \) es la densidad del fluido.

- \( v \) es la velocidad del fluido.

- \( g \) es la aceleración debido a la gravedad.

- \( h \) es la altura del fluido respecto a un punto de referencia.
}

\info{Simplificación de Bernoulli}{
si dividimos ambos miembros por el peso específico
\begin{equation}
    \dfrac{p}{\gamma}+\dfrac{v^2}{2g}+h=constante
\end{equation}

siendo:

$h:$ Altura del líquido.[m]

$\gamma:$ Peso específico del líquido [$N/m^3$]

$g:$ Aceleración de la gravedad $g=9,8 m/s^2$

$v:$ Velocidad del fluido [$m/s$]
}



\section{Teorema de Torricelli}

\info{Definición de Teorema de Torricelli}{
La velocidad de salida de un líquido a través de un pequeño orificio es la misma que adquiriría un objeto que cayera libremente en el vacío dese una altura $H$

\begin{equation}
    v=\sqrt{2gH}
\end{equation}

siendo:

$H:$ Altura del líquido [m]

$g:$ Aceleración de la gravedad $m/s^2$

$v:$ Velocidad de salida del líquido [$m/s$]
}

\actividad{Resolver los siguientes problemas.}
\begin{enumerate}
    \item Un tanque cilíndrico de de $r=50cm$ de radio contiene una masa de agua de $m=10000kg$, si se perfora un pequeño orificio por la parte baja, encontrar la velocidad $v_2$ de salida.
    \item ¿Con qué velocidad sale un líquido por un orificio que se encuentra a una profundidad de 3 m?
    \item  Determine a qué altura se debe abrir un orificio de un estanque, para que el líquido salga con una velocidad de $v=9,8 m/s$
    \item Determinar la magnitud de la velocidad con la que sale un líquido por un orificio localizado a una profundidad de $H=2,6m$ en un tanque de almacenamiento.
    \item  Determinar a qué altura se debe abrir un orificio de un tanque de nafta, para que el líquido salga con una velocidad de $v=13,5 m/s$. 
    \item Para medir la magnitud de la velocidad de la corriente en un rió se introduce un tubo de Pitot, la altura a la que llega el agua dentro del tubo es de $H=0,20m$ ¿a qué magnitud de velocidad va la corriente?
    \item Se desea que un estanque posea circulación de agua, para ello se practica un orificio a de $r=5mm$ bajo $H=0.5m$ del nivel del agua determinar el caudal necesario de recirculación $Q=S \cdot v$.
    \item  Un flujo de agua va de la $s_1$ a la $s_2$. La $s_1$ tiene $d_1=25mm$, la presión manométrica es de $p_1=345kPa$, y la velocidad de flujo es de $v_1=3 m/s$. La $s_2$, mide $d_2=50mm$ de diámetro, y se encuentra a 2 metros por arriba de la $s_1$. Si suponemos que no hay pérdida de energía en el sistema. Calcular la presión $p_2$.
    \item Un tanque de sección cuadrada de lado $l = 50 cm$ contiene una masa de agua de $m = 31000kg$. Si se perfora un pequeño orificio en la parte inferior, encuentra la velocidad $v_2$ de salida.

\item ¿Con qué velocidad sale un líquido por un orificio de una pileta de natación que se encuentra a una profundidad de $30m$?

\img{capitulos/images/06/pileta.jpg}{Pileta de natación}

\item Determina a qué altura se debe abrir un orificio en un estanque para que el líquido salga con una velocidad de $v = 1000 , \text{m/s}$.

\item Determina la magnitud de la velocidad con la que sale un líquido por un orificio localizado a una profundidad de $H = 30 , \text{m}$ en un tanque de almacenamiento.

\item Determina a qué altura se debe abrir un orificio en un tanque de aceite para que el líquido salga con una velocidad de $v = 15 \text{m/s}$.

\item Para medir la magnitud de la velocidad de la corriente en un río se introduce un tubo de Pitot. La altura a la que llega el agua dentro del tubo es de $h=0,20m$. ver video.

\video{https://www.youtube.com/watch?v=AQviCLqrznc}{Tubo de pitot casero.}




\item Un recipiente tiene forma de cono invertido y contiene agua hasta una altura de $h=8m$. Si se practica un orificio en la parte inferior, determine la velocidad de salida del agua en función del tiempo.

\item En un tanque con forma de prisma triangular, el área de la base es de $A = 10 \, \text{m}^2$ y la altura es de $h = 5 \, \text{m}$. Si se abre un orificio en la base, calcule la tasa de cambio de la velocidad de salida del agua en función del tiempo.

\item Un tanque biodigestor de producción de biogas está lleno de aguas negras hasta una altura de $h = 12 \, \text{m}$. Si se abre un orificio en la parte inferior, determine  la velocidad de salida.

\img{capitulos/images/06/biodigestor.jpg}{Biodigestor para producción de biogas.}

\item Un tanque tiene forma de esfera y contiene agua hasta una altura de $h = 4 \, \text{m}$. Si se practica un orificio en la parte inferior, determine la ecuación que relaciona la velocidad de salida del agua con el radio del orificio y el tiempo.

\item En un depósito con forma de prisma rectangular, el ancho es de $w = 6 \, \text{m}$, el largo es de $l = 8 \, \text{m}$ y la altura es de $h = 10 \, \text{m}$. Si se abre un orificio en la base,de 1cm de diámetro determine la velocidad con que baja la altura del líquido del tanque apenas comienza a descargarse.

\item Un recipiente de forma irregular contiene agua hasta una altura de $h = 15 \, \text{m}$. Si se practica un orificio pequeño, determine la velocidad de salida del líquido.




\end{enumerate}

\end{document}